El desarrollo de este trabajo surge con el objetivo de añadir funcionalidades que no están presentes en otras plataformas existentes. 
Hasta ahora, los usuarios solo podían valorar videojuegos, almacenarlos en listas personalizadas y buscarlos por género, título o año 
de lanzamiento. Sin embargo, estas aplicaciones no permiten una interacción real entre usuarios, lo que limita su potencial como red 
social centrada en los videojuegos.

Además, algunas necesidades específicas, como la búsqueda de videojuegos según la desarrolladora, no están contempladas en estas 
plataformas. Por todo ello, este proyecto propone el desarrollo de una aplicación web que permita una mayor interacción entre usuarios 
y amplíe las posibilidades de búsqueda, ofreciendo así una experiencia más completa y personalizada.